\begin{document}
\maketitle

\begin{abstract}
This paper proves the algebraic core of TFPT's carrier theorem and, separately, sketches its
Standard-Model packet outlook. The carrier core takes four explicit interfaces --- a carrier
involution $\epscar^2 = \mathbf 1$, a Higgs rank-$2$ interface $\dim E_+ = 2$, a Yukawa Stage-D
trilinear interface $\omega : \Lambda^3 E_- \to K$ with operational contraction predicates, and a
primitive oriented determinant-preserving carrier condition --- and forces the determinant-normalised
$3+2$ hypercharge carrier
\[
Y = -\tfrac13 P_- + \tfrac12 P_+ = \operatorname{diag}\!\left(-\tfrac13,-\tfrac13,-\tfrac13,\tfrac12,\tfrac12\right),
\quad
(q_-, q_+) = (-2, 3),
\quad
\tr_E Y = 0,
\quad
6 Y^2 - Y - \mathbf 1 = \mathbf 0.
\]
The algebraic dependency graph is \emph{formally verified} in Lean 4 (companion note;
\texttt{audit.sh} returns \texttt{AUDIT:~PASS}). The remaining downstream content of TFPT's
Standard-Model packet --- the even exterior spinor packet $S^+ = \Lambda^{\mathrm{even}} E$,
derived representations (duals, exterior powers, tensor products), the $\mathbb Z_6$ quotient
$G_{\mathrm{phys}}$, the family count $N_{\mathrm{fam}} = 3$, the admissible occupancy
$\Omega_{\mathrm{adm}} = 48$, the compact Higgs index $N_\Phi = 1$, and the abelian index
coefficient $b_1 = 41/10$ --- is presented as a \emph{conditional outlook} downstream of the
carrier core; not as a formalised theorem.
\end{abstract}

\begin{tfptfrontbox}
\textbf{Inputs from previous papers.} The primitive boundary kernel
$(\taudbl,\iC,\Pprim,[u_\Sigma],\cthree)$ from Paper 1, restricted to the finite essential carrier $E$
as an inhabitant of \texttt{BoundaryPolarization}.

\textbf{New theorem contribution.}\ The $3+2$ hypercharge carrier rigidity theorem, as a
Lean-verified algebraic statement, plus the Lean-typed targets for the upstream and downstream
content of Paper 2.

\textbf{Not claimed here.} No exact $\alpha$ value, no CKM/PMNS closure, no pole-mass ledger, no OS
reconstruction, no cosmology, no E8 grammar; and (newly explicit) \emph{not} the full SM hypercharge
spectrum on derived representations, \emph{not} the $\mathbb Z_6$ quotient as a Lean theorem,
\emph{not} the family count or admissible occupancy as Lean theorems --- these are downstream
Status-C readouts.

\textbf{Falsification or audit surface.} The paper fails if the carrier normal form imports
Standard-Model representation data by hand, if the Higgs / Stage-D-Yukawa discharge is not made
explicit, or if downstream Status-C readouts are presented as Status-A theorems.
\end{tfptfrontbox}


\begin{tfptclaimcontract}
\textbf{Top-level theorem.}\ Given the four explicit interfaces below, the determinant-normalised
$3+2$ hypercharge carrier is uniquely
$Y = \operatorname{diag}(-\tfrac13,-\tfrac13,-\tfrac13,\tfrac12,\tfrac12)$, with $(q_-,q_+)=(-2,3)$,
$\tr_E Y = 0$, and $6 Y^2 - Y - \mathbf 1 = \mathbf 0$ in $\operatorname{End}_K E$
\emph{(Status A: Lean theorem
\texttt{CarrierPremises.hypercharge\_carrier\_packet})}.

\smallskip
\textbf{Direct inputs of the Lean theorem.}
\begin{enumerate}[leftmargin=1.7em, itemsep=0pt, label=(\arabic*)]
    \item carrier involution $\epscar^2 = \mathbf 1$ on $E$
        \emph{(Status B: \texttt{CarrierInvolution} / \texttt{BoundaryPolarization})};
    \item Higgs rank-$2$ interface $\dim E_+ = 2$
        \emph{(Status B: \texttt{HiggsIndexCertificate}, refined by
        \texttt{HiggsTopForm} and \texttt{HiggsSchemeCohomologyShadow})};
    \item Yukawa Stage-D interface: an alternating trilinear form
        $\omega : \Lambda^3 E_- \to K$ with
        \texttt{ContractionInjective}, \texttt{ContractionSurjective},
        and $\texttt{Nontrivial}\,E_-$
        \emph{(Status A: $\dim E_- = 3 \iff$ such an $\omega$ exists; see
        \texttt{YukawaStageDExistence.stage\_D\_iff\_finrank\_eq\_three})};
    \item primitive oriented determinant-preserving carrier condition
        $3 q_- + 2 q_+ = 0,\ q_- < 0 < q_+,\ \gcd(|q_-|, q_+) = 1$
        \emph{(Status B: \texttt{PrimitiveOrientedDeterminantCarrier},
        re-named via \texttt{SeamWindingData})}.
\end{enumerate}

\smallskip
\textbf{Open upstream TFPT-origin theorems.}\ Each direct input has a Lean-typed target and an
open TFPT-side derivation:
\begin{enumerate}[leftmargin=1.7em, itemsep=0pt, label=(\roman*)]
    \item Paper 1 analytic boundary theory $\Rightarrow$ \texttt{BoundaryPolarization}
        (analytic Calder\'on construction);
    \item TFPT boundary kernel + representation data $\Rightarrow$
        \texttt{BoundaryYukawaKernel} (canonical $\omega$);
    \item seam winding $[u_\Sigma] = 1$ $\Rightarrow$ \texttt{SeamWindingData}
        (orientation, primitivity, trace-zero).
\end{enumerate}

\smallskip
\textbf{Downstream Status-C outlook.}\ Conditional on the carrier core, this paper \emph{outlines}
(does not formalise) the derived SM representations, the $\mathbb Z_6$ quotient
$G_{\mathrm{phys}}$, the one-family spinor packet $S^+ = \Lambda^{\mathrm{even}} E$, the family
count $N_{\mathrm{fam}} = 3$, the admissible occupancy $\Omega_{\mathrm{adm}} = 48$, and the
abelian index coefficient $b_1 = 41/10$.

\smallskip
\textbf{Kill condition.}\ The Lean-verified Status-A core fails iff any check in the companion
audit (build / no-sorry / no-axiom / no-unsafe / axiom-print / \texttt{\#check} headlines /
signature-lock examples) fails. The Status-C downstream content is auditable through its
explicit downstream-conditionality rather than by Lean verification.
\end{tfptclaimcontract}

\begin{tfptstatusbox}
\textbf{\stag{L} (Lean verified, this paper's core theorem).}\
The hypercharge carrier packet
$\bigl((q_-,q_+)=(-2,3),\ \tr_E Y = 0,\ 6 Y^2 - Y - \mathbf 1 = \mathbf 0\bigr)$
is a Lean 4 theorem
(\texttt{CarrierPremises.hypercharge\_carrier\_packet}),
proved from the four direct inputs of the claim contract,
with all upstream refinements (\texttt{BoundaryPolarization},
\texttt{HiggsTopForm}, \texttt{YukawaTrilinearForm} +
\texttt{YukawaStageDExistence}, \texttt{SeamWindingData})
Lean-typed and Lean-bridged. The signature-locked
\texttt{AuditContract.lean} guards every headline.

\smallskip
\textbf{\stag{I} (Lean interfaced, open TFPT-origin theorems).}\
Three typed Lean targets carry the upstream TFPT obligations:
\texttt{BoundaryPolarization} (analytic Paper-1 content),
\texttt{BoundaryYukawaKernel} (TFPT-canonical $\omega$),
\texttt{SeamWindingData} (seam-class topological content).
The Higgs Stage-B algebraic shadow is \stag{L} but the
$\mathrm{Proj}$-side identification
$H^0(\mathrm{Proj}\,K[X,Y], \mathcal O(1)) \simeq K[X,Y]_1$
is \stag{S} (standard algebraic-geometry theorem cited from
the literature) with a manuscript discharge sketched in
Paper~5.

\smallskip
\textbf{\stag{M} (manuscript theorems in this paper's body).}\
The source-extract body contains \stag{M} theorems on the
internal carrier stabilizer, the faithful physical global
gauge group construction (input to the $\mathbb Z_6$ quotient),
the one-family exterior packet, the manuscript-level
"carrier rigidity from the primitive component of the branch
Yukawa tensor", and the joint discrete admissible sector.
These are \emph{not} \stag{L}.

\smallskip
\textbf{\stag{C} (mathematical proposal, downstream outlook).}\
Derived SM representations, $\mathbb Z_6$ quotient as a Lean
theorem, family count, admissible occupancy, abelian index,
and observable readouts are presented as conditional
consequences of the \stag{L} core and the \stag{M} body,
not as formalised theorems.

\smallskip
The companion Lean note is
\texttt{experiments/lean4-carrier-rigidity/note\_carrier\_rigidity\_lean4.tex}
and the reproducible build is
\texttt{experiments/lean4-carrier-rigidity/scripts/audit.sh}.
\end{tfptstatusbox}

\tableofcontents

\section{Formal Status of the Carrier Theorem}
\label{sec:formal-status}

This paper separates the \emph{carrier rigidity theorem}
(Status A, Lean verified) from the \emph{Standard-Model packet
outlook} (Status C, downstream).

The Status-A theorem is the algebraic dependency graph
\[
\big(\epscar,\ \dim E_+ = 2,\ \omega \text{ with CI, CS, } \texttt{Nontrivial}\,E_-,\ \text{primitive oriented det.\,carrier}\big)
\;\;\Longrightarrow\;\;
\big(Y,\ (q_-, q_+),\ \tr Y,\ 6 Y^2 - Y - \mathbf 1\big).
\]
This is the statement we prove formally and audit. The
input objects on the left are Status-B Lean-typed targets;
the output is a Status-A Lean theorem.

The Status-C outlook --- the even exterior spinor packet
$S^+ = \Lambda^{\mathrm{even}} E$, derived SM representations,
the $\mathbb Z_6$ quotient $G_{\mathrm{phys}}$, the family
count, admissible occupancy, and the abelian index --- is a
\emph{conditional} consequence of the Status-A theorem. It is
explicitly \emph{not} on the same formalisation level. We
include it because TFPT predicts it; we mark it as Status C
because no Lean proof exists.

\noindent
\textbf{Lean verification summary.}
The seven-check audit of the companion Lean note returns
\texttt{AUDIT:~PASS}: \texttt{lake build} succeeds; no
\texttt{sorry} or \texttt{admit}; no axiom or constant
declarations; no \texttt{unsafe}/\texttt{opaque}/\texttt{partial}
escape hatches; every headline theorem depends only on
$\{\texttt{propext},\ \texttt{Classical.choice},\ \texttt{Quot.sound}\}$;
every headline elaborates via Lean \texttt{\#check}; every
headline type is signature-locked via an
\texttt{example : <exact intended type> := <theorem>} pattern.

\section{Carrier Polynomial as a Corollary}

The paper does not begin by assuming
\[
6Y^2-Y-\mathbf 1=0,\qquad \tr_EY=0.
\]
It begins with the carrier involution induced by the boundary polarization:
\[
\varepsilon_{\mathrm{car}}:=\iC|_E,
\qquad
\varepsilon_{\mathrm{car}}^2=\mathbf 1,
\qquad
E=E_-\oplus E_+,
\qquad
P_\pm=\frac{\mathbf 1\pm\varepsilon_{\mathrm{car}}}{2}.
\]
At this stage there is only a two-point algebra
\[
\mathbb C[\varepsilon_{\mathrm{car}}]=\{a\mathbf 1+b\varepsilon_{\mathrm{car}}\}.
\]
No $3+2$ split, hypercharge vector, Standard-Model packet, or electromagnetic number has been
used.

\begin{theorem}[$3+2$ Hypercharge Carrier Rigidity --- Lean verified]
\label{thm:carrier-rigidity}
Let $K$ be a field of characteristic zero and let $E$ be a finite-dimensional $K$-vector space
equipped with a carrier involution $\epscar : E \to E$ with $\epscar^2 = \mathrm{id}_E$. Set
$P_\pm = (1 \pm \epscar)/2$ and $E_\pm := P_\pm E$. Assume:
\begin{enumerate}[label=(\roman*), leftmargin=1.7em]
    \item \emph{(Higgs Stage B / Stage B+).}\ A Higgs Stage-B
        certificate $E_+ \simeq_K \texttt{DegreeOneTwoVarForm}\,K$
        (equivalently, $E_+ \simeq_K K[X_0, X_1]_1$ via the
        Lean-verified
        $\texttt{HiggsSchemeCohomologyShadow.toDegreeOneTwoVarForm}$),
        forcing $\dim E_+ = 2$.
    \item \emph{(Yukawa Stage D).}\ A non-trivial finite-dimensional
        Stage-D datum: an alternating trilinear form
        $\omega : \Lambda^3 E_- \to_{\ell K} K$ whose associated
        contraction
        $C_\omega(v)(w \wedge x) := \omega(v \wedge w \wedge x)$ is
        injective \emph{(\texttt{ContractionInjective})} and
        surjective \emph{(\texttt{ContractionSurjective})}, with
        $\texttt{Nontrivial}\,E_-$.
        Equivalently (Lean-verified characterisation), $\dim E_- = 3$.
    \item \emph{(Primitive oriented determinant-preserving carrier
        condition.)}\ Integer charges $q_-, q_+ \in \mathbb Z$ with
        \[
        3 q_- + 2 q_+ = 0,
        \qquad
        q_- < 0 < q_+,
        \qquad
        \gcd(|q_-|, q_+) = 1.
        \]
\end{enumerate}
Then $\dim E_+ = 2$, $\dim E_- = 3$, $(q_-, q_+) = (-2, 3)$, and the
determinant-normalised two-point generator is uniquely
\[
Y = -\tfrac13 P_- + \tfrac12 P_+
  = \operatorname{diag}\!\left(-\tfrac13,-\tfrac13,-\tfrac13,\tfrac12,\tfrac12\right).
\]
Furthermore
\[
\tr_E Y = 0,
\qquad
6 Y^2 - Y - \mathbf 1 = \mathbf 0
\quad\text{in}\quad
\mathrm{End}_K E.
\]
\end{theorem}

\begin{proof}[Proof spine, with Lean-verified steps]
The theorem is the bundled top-level theorem
\texttt{CarrierPremises.hypercharge\_carrier\_packet} of the
companion Lean note. We record the spine:

\begin{enumerate}[label=(S\arabic*), leftmargin=2em, itemsep=2pt]
    \item \emph{Boundary involution $\Rightarrow$ projectors.}\
        $\epscar^2 = \mathbf 1$ together with
        $\mathrm{char}\,K \ne 2$ gives complete orthogonal
        idempotents $P_-, P_+$ with $P_- + P_+ = \mathbf 1$
        (\texttt{InvolutionProjectors.CarrierInvolution.toPolarization}).
    \item \emph{Higgs Stage B $\Rightarrow$ $\dim E_+ = 2$.}\ The
        algebraic shadow of $H^0(\mathbb P^1, \mathcal O(1))$ has
        rank $2$ over $K$ (\texttt{HiggsTopForm.finrank\_Eplus\_eq\_two}).
    \item \emph{Yukawa Stage D $\Rightarrow$ $\dim E_- = 3$.}\ With
        $\texttt{Nontrivial}\,E_-$ and both contraction predicates,
        the contraction $C_\omega$ is a linear isomorphism
        $E_- \simeq (\Lambda^2 E_-)^*$ on a finite-dimensional space,
        so the dimension chain
        $\dim E_- = \dim(\Lambda^2 E_-)^* = \binom{\dim E_-}{2}$
        and elementary integer arithmetic force
        $\dim E_- = 3$
        (\texttt{YukawaTrilinearForm.finrank\_E\_eq\_three}).
        The backward direction --- that any such $E_-$ of dimension
        $3$ admits such an $\omega$, constructed from a basis via
        $b.\texttt{det}$ --- is
        \texttt{YukawaStageDExistence.exists\_yukawaTrilinearForm\_of\_finrank\_eq\_three};
        the joint characterisation under $\texttt{Nontrivial}$ is
        \texttt{stage\_D\_iff\_finrank\_eq\_three}.
    \item \emph{Primitive oriented determinant-preserving condition
        $\Rightarrow$ $(q_-, q_+) = (-2, 3)$.}\ For $(m, n) = (3, 2)$,
        the integer Diophantine equation
        $3 q_- + 2 q_+ = 0$ together with the orientation and
        primitivity conditions has the unique solution
        $(-2, 3)$
        (\texttt{LatticeRigidityGeneral.primitive\_trace\_free\_pair\_general}
        specialised to $(3, 2)$;
        \texttt{Rigidity.unique\_carrier\_pair}).
    \item \emph{Trace identity.}\ With $a = -\tfrac13$, $b = \tfrac12$,
        $\dim E_- = 3$, $\dim E_+ = 2$, the structural trace formula
        $\tr_E (a P_- + b P_+) = a \dim E_- + b \dim E_+$ gives
        $\tr_E Y = -1 + 1 = 0$
        (\texttt{TraceProjection.trace\_carrier\_Y\_eq\_zero}).
    \item \emph{Carrier polynomial.}\ The minimal polynomial of $Y$
        with roots $-\tfrac13, \tfrac12$ is
        $\left(Y + \tfrac13\right)\left(Y - \tfrac12\right) = 0$,
        equivalently $6 Y^2 - Y - \mathbf 1 = \mathbf 0$
        in $\mathrm{End}_K E$
        (\texttt{CarrierPremises.carrier\_polynomial\_Y}).
\end{enumerate}

Steps (S1)--(S6) are individually Lean theorems and are
assembled into \texttt{CarrierPremises.hypercharge\_carrier\_packet}.
Every signature is locked in
\texttt{AuditContract.lean}.\qedhere
\end{proof}

\begin{remark}[Where the language "primitive indecomposable Yukawa type" sits]
\label{rem:primitive-yukawa-naming}
In earlier presentations the rank-$3$ discharge was stated as
``primitive indecomposable Yukawa type'' without an algebraic
definition. The Lean formalisation forced a discipline: the
operational content of the phrase is precisely the conjunction
\[
\texttt{ContractionInjective}\ \omega
\;\wedge\;
\texttt{ContractionSurjective}\ \omega
\;\wedge\;
\texttt{Nontrivial}\,E_-,
\]
where ``\emph{non-degenerate}'' is the docstring synonym of
\texttt{ContractionInjective} and ``\emph{primitive / no
spectator}'' is the docstring synonym of
\texttt{ContractionSurjective}. The phrase
``primitive indecomposable Yukawa type'' is retained as
mathematical prose; it no longer enters the formal hypothesis
of \Cref{thm:carrier-rigidity}.

The Lean theorem
\texttt{YukawaTrilinearForm.finrank\_E\_eq\_zero\_or\_three}
records the transparency form: without
$\texttt{Nontrivial}\,E_-$, the operational predicates yield
$\dim E_- = 0 \vee \dim E_- = 3$; non-triviality is the
algebraic content of "indecomposable" that rules out the
trivial case.
\end{remark}

\begin{remark}[General split check]
For a general essential split $E_b\oplus E_s$ with determinant-normalized roots
$-1/b$ and $1/s$, the two-point generator satisfies
\[
bs\,Y^2+(s-b)Y-\mathbf 1=0.
\]
The Higgs Stage-B and Yukawa Stage-D arguments fix $(b,s)=(3,2)$, hence the coefficient
$6=3\cdot2$ is not guessed or phenomenologically tuned.
\end{remark}

